\chapter{Revisão da Literatura} \label{sec:revisao}

\section{Métodos de Modelagem de Risco de Crédito}

A literatura sobre modelagem de risco de crédito evoluiu substancialmente ao longo das últimas décadas. O trabalho seminal de \cite{hand&henley1997} revisou sistematicamente os métodos de classificação aplicados ao \textit{credit scoring}, comparando regressão logística, análise discriminante e k-vizinhos mais próximos (\textit{k-Nearest Neighbors} — k-NN), concluindo que a regressão logística oferece um melhor equilíbrio entre poder preditivo e interpretabilidade. \cite{thomasetal2002} consolidaram em obra de referência os métodos quantitativos para \textit{credit scoring} aplicáveis a portfólios de varejo e empresarial, abordando tanto a modelagem estatística quanto os aspectos de implementação operacional.

Com o avanço do aprendizado de máquina, os métodos \textit{ensemble}, ganharam poeminência nas aplicações de crédito. \cite{breiman2001} propôs o algoritmo \textit{Random Forest}, que constrói múltiplas árvores de decisão descorrelacionadas e agrega suas previsões por votação ou média, demonstrando alta robustez contra overfitting e elevada capacidade preditiva em datasets de alta dimensionalidade. \cite{friedman2001} desenvolveu o framework de \textit{Gradient Boosting}, que se tornou a base para implementações de grande impacto como o \textit{XGBoost} \cite{chen&guestrin2016}, amplamente utilizado em competições de modelagem de crédito e na indústria financeira global por sua escalabilidade e precisão.

No contexto específico do crédito rural, \cite{turvey1991} explorou a aplicação de \textit{credit scoring} a empréstimos agropecuários, identificando os desafios particulares desse segmento: variabilidade climática, dependência de ciclos de preços de commodities e a natureza cíclica do processo produtivo agrícola. \cite{barry&ellinger2011} discutem como as finanças agrícolas lidam com estruturas de risco distintas do crédito comercial convencional, exigindo variáveis específicas vinculadas à produção agrícola, ciclo operavional , condições de mercado e fatores climáticos.

Mais recentemente, a literatura tem abordado o desafio da interpretabilidade em modelos de aprendizado de máquina aplicados ao risco de crédito. \cite{lundberg&lee2017} propuseram o \textit{framework SHAP} (\textit{SHapley Additive exPlanations}), fundamentado na teoria dos jogos cooperativos, para interpretar as previsões de modelos complexos por meio da decomposição aditiva da contribuição de cada variável. O SHAP tem sido crescentemente adotado no setor financeiro para conciliar o poder preditivo de modelos complexos com a necessidade regulatória de explicabilidade das decisões de crédito.

\section{Dados Abertos e Regulatórios para Crédito Rural}

O uso de dados abertos e regulatórios para a análise de risco de crédito rural e a previsão de inadimplência é uma área em expansão tanto na literatura acadêmica quanto nas práticas de mercado. Em âmbito internacional, estudos têm explorado o uso de dados de sistemas oficiais de seguro agrícola, registro de crédito regulatórios e variáveis macroeconômicas para construir modelos que estimam o risco sem a necessidade de dados comportamentais próprios dos tomadores \cite{barry&ellinger2011}. Essa abordagem é especialmente relevante para novos entrantes e \textit{fintechs} de crédito rural, que carecem de bases históricas internas suficientes para a calibração de modelos convencionais de \textit{credit scoring}.

No contexto brasileiro, os microdados do SICOR representam uma fonte singular para esse tipo de análise. O Bacen, por meio de seus Anuários Estatísticos, documenta sistematicamente o panorama estatístico de originação, inadimplência e distribuição do crédito rural por cultura e região, estabelecendo uma base de referência empírica para o setor. Estudos baseados nessas estatísticas têm identificado padrões de inadimplência diferenciada entre culturas, regiões e safras, apontando para a forte influência de eventos climáticos adverosos (secas, geadas, etc.) e de ciclos de preços de \textit{commodities} nos índices de \textit{default}.

A integração de variáveis setoriais — como atividade, finalidade do financiamento e porte do tomador — com variáveis macroeconômicas — como câmbio, taxa Selic e Plano Safra, preços de \textit{commodities} e índices de produção agrícola — tem demostrado potencial para aumentar substancialmente a capacidade preditiva de modelos agnósticos. Essa abordagem encontra respaldo na literatura de modelagem de portfólio de crédito, que reconhece a importância dos fatores sistémicos e setoriais na determinação da probabilidade de \textit{default} \cite{thomasetal2002}. A disponibilidade pública e a cobertura temporal do SICOR posicionam esse dados como um ativo analítico de alto valor para estratégias \textit{data-driven} de entrada no mercado de crédito rural.
